% Reunión expositiva de precedencia causal entre resultados ya establecidos.
% Los propietarios íntegros continúan gobernando cada flecha.
% Propietarios reunidos por el grafo activo:
% - sucesor_102/deltas_ley9/c27_teorema_generacion_coinductiva_profundidad_arbitraria.tex
% - sucesor_102/deltas_ley9/c34_diagrama_bifurcado_y_decada_accion.tex
% - sucesor_102/deltas_ley9/c36_recta_accion_planck.tex
% - colaboracion/parte_iii/tex/capitulos_finales/32_raiz_menos_uno.tex
% - incorporaciones/parte_iv_ampliaciones_20260824/sources/editor_ready/03_pareja_planck_respuesta.tex
% - integracion_83/parte_iv/chapters/48_electron_primera_realizacion.tex

\section{Généalogie opératoire de l'action, de la géométrie angulaire et de la masse}
\label{sec:l9s102:transicion-constantes-accion-geometria-masas}

La dynamique coinductive APP--TRIT--TPK se prolonge à profondeur arbitraire
sur un même état enrichi et sur la structure discrète conjointe du
continu. À chaque prolongation, elle conserve valeur, généalogie, phase, retenue,
orientation, vacance, chemin, frontière, mémoire et survie. Ses
publications distinguées sont d'abord les trois coordonnées internes et,
par le sceau dodécaphasique, la coordonnée \(\alpha_{\mathrm{HMT}}\) :
\[
 \mathrm{APP}\longrightarrow\mathrm{TRIT}\longrightarrow\mathrm{TPK}
 \longrightarrow\mathcal S_{\infty}^{\mathrm{enr}}
 \longrightarrow\mathcal C_{\mathrm{cont}}^{\mathrm{disc}}
 \longrightarrow(\pi_{\mathrm{HMT}},\varphi_{\mathrm{HMT}},e_{\mathrm{HMT}})
 \longrightarrow\alpha_{\mathrm{HMT}}.
\]
La dernière flèche représente la publication signée de
\(\alpha_{\mathrm{HMT}}\) au sein de cette même dynamique ; elle ne reçoit pas de valeurs
conventionnelles et ne transforme pas les \(\pi,\varphi,e\) conventionnels en générateurs.

Ce n'est qu'ensuite que l'on applique le lecteur déterminantal d'action. Celui-ci publie la
valuation \(10^{-34}\mathcal U_S\), les deux sections
\(\hbar_{\mathrm{pre}},\hbar_{\mathrm{ret}}\) et leurs actions de cycle
complet
\[
 h_{\mathrm{pre}}=2\pi_{\mathrm{HMT}}\hbar_{\mathrm{pre}},
 \qquad
 h_{\mathrm{ret}}=2\pi_{\mathrm{HMT}}\hbar_{\mathrm{ret}}.
\]
À partir de ce point causal, deux branches restent séparées et ne se
sélectionnent pas mutuellement. La branche dimensionnelle--gravitationnelle prolonge l'action et
la section causale jusqu'aux formes déjà démontrées du couple de Planck,
\[
 \ell^{2}=G\hbar c^{-3},
 \qquad
 \frac{\ell}{Q}=Gc^{-4};
\]
la branche complexe publie, à partir de l'opérateur interne de rotation,
\[
 u^2+1=0,
 \qquad u\longmapsto J_{+1},
 \qquad J_{+1}^{2}=-I.
\]
Telle est la réalisation opératoire interne de \(i\), non l'introduction d'un
symbole complexe sans antécédent. Les branches dimensionnelle--gravitationnelle et
complexe sont parallèles sur l'état enrichi commun : la première ne
transforme pas \(i\) en entrée de l'échelle d'action et, réciproquement,
\(h\), \(\hbar\), \(c^{-3}\) et \(c^{-4}\) ne génèrent pas \(i\). Les identités
du couple de Planck restent ainsi séparées de la réalisation complexe.

En écrivant \(\prec_{\mathrm{HMT}}\) pour la précédence causale avec conservation
du type ---et non pour affirmer que chaque objet de droite est fonction de
tous ceux de gauche---, on obtient
\[
\begin{aligned}
 (\pi,\varphi,e,\alpha)_{\mathrm{HMT}}
 &\prec_{\mathrm{HMT}}(10^{-34}\mathcal U_S,\hbar,h)\\
 &\prec_{\mathrm{HMT}}
 \bigl\{G\hbar c^{-3},Gc^{-4};\,J_{+1}^{2}=-I, i\bigr\}\\
 &\prec_{\mathrm{HMT}}(A_{\mathrm{deg}},C^{*}_{\mathrm{deg}})\\
 &\prec_{\mathrm{HMT}}(\text{électron},\text{loi des masses}).
\end{aligned}
\]
La construction angulaire conserve ses entrées effectives
\(\alpha_{\mathrm{HMT}}\) et \(\eta_{\mathrm{ret}}\) : les deux branches parallèles
précédentes fixent la disponibilité conjointe et le typage, mais ne sont pas promues
rétrospectivement au rang de sélecteurs de l'angle ou du contre-angle.

Dans le chemin électronique central, les lecteurs bidirectionnels conservent deux
étapes distinctes d'une seule généalogie :
\[
 \begin{aligned}
 E_e^{(0)}
   &=0.5109989224880660725\ldots\ \mathrm{MeV},
   &&\text{étape basale},\\
 E_e^{\beta}
   &=E_e^{(0)}R_{\mathrm{act}}^{\,80-54\alpha_{\mathrm{HMT}}+6\Delta_4}
     =0.510998950690000808608\ldots\ \mathrm{MeV},
   &&\text{sortie bêta directrice},
 \end{aligned}
\]
avec
\[
 K_D(\Gamma_e)=K_{\Omega}(\Gamma_e)=(80,54,6).
\]
La coïncidence du triplet dans la fibre centrale n'identifie ni les domaines ni
les opérations de \(K_D\) et \(K_{\Omega}\). Dans le registre canonique
historique, \texttt{MD-ELECTRON-001} désigne le candidat étalonné de l'étape
basale. Cette identification conserve sa fonction de provenance, mais ne gouverne pas
la sortie bêta et ne sélectionne pas le chemin ; le basal et la clôture bêta restent
des étapes typées des lecteurs bidirectionnels.

\paragraph{Critère de réfutation.}
La généalogie précédente est réfutée si une valeur conventionnelle ou la sortie
bêta est utilisée pour sélectionner un chemin, si l'étape basale est présentée comme
clôture directrice, si l'une des deux branches parallèles est identifiée sans une
application typée, ou si l'angle, le contre-angle, l'électron ou une masse sont placés
comme entrées du générateur de constantes ou d'action.
